%==============================================================
\section{李群的基本概念}
\label{sec:lie-group}
%==============================================================

\subsection{定义}

\textbf{李群} $G$ 是一个同时具有群结构和光滑流形结构的对象，使得群乘法 $G \times G \to G$ 和取逆 $G \to G$ 都是光滑映射。

\textbf{维度：} 李群作为流形的维度。

\subsection{经典李群}

\begin{center}
\begin{tabular}{p{2cm}p{5cm}p{1.5cm}p{4cm}}
\toprule
\textbf{群} & \textbf{定义} & \textbf{维度} & \textbf{应用} \\
\midrule
$GL(n,\mathbb{R})$ & $n\times n$ 可逆实矩阵 & $n^2$ & 一般线性变换 \\
$SL(n,\mathbb{R})$ & $\det A = 1$ & $n^2 - 1$ & 保体积变换 \\
$O(n)$ & $A^TA = I$ & $\frac{n(n-1)}{2}$ & 保距（含反射） \\
$SO(n)$ & $A^TA = I$，$\det A = 1$ & $\frac{n(n-1)}{2}$ & 旋转 \\
$U(n)$ & $A^\dagger A = I$ & $n^2$ & 量子力学 \\
$SU(n)$ & $A^\dagger A = I$，$\det A = 1$ & $n^2 - 1$ & 规范场论 \\
$Sp(n)$ & $A^\dagger A = I$，保辛形式 & $n(2n+1)$ & 哈密顿力学 \\
$SE(n)$ & $\begin{pmatrix} R & t \\ 0 & 1 \end{pmatrix}$，$R \in SO(n)$ & $\frac{n(n+1)}{2}$ & 刚体运动 \\
\bottomrule
\end{tabular}
\end{center}

\subsection{关键低维例子}

\subsubsection{$SO(2)$：平面旋转}

\[
SO(2) = \left\{\begin{pmatrix} \cos\theta & -\sin\theta \\ \sin\theta & \cos\theta \end{pmatrix} : \theta \in \mathbb{R}\right\}
\]

\begin{itemize}
  \item 维度：1（一个参数 $\theta$）
  \item 拓扑：$S^1$（圆）
  \item 阿贝尔群（交换）
  \item 同构于 $U(1)$（复数单位圆）
\end{itemize}

\subsubsection{$SO(3)$：三维旋转}

\[
SO(3) = \{R \in \mathbb{R}^{3\times 3} : R^TR = I,\; \det R = 1\}
\]

\begin{itemize}
  \item 维度：3（三个参数，如欧拉角或轴角）
  \item 拓扑：$\mathbb{RP}^3$（三维实射影空间）
  \item 非阿贝尔群
  \item 基本群：$\pi_1(SO(3)) = \mathbb{Z}_2$（旋转 $2\pi$ 不等于恒等！）
\end{itemize}

\subsubsection{$SU(2)$：自旋群}

\[
SU(2) = \{A \in \mathbb{C}^{2\times 2} : A^\dagger A = I,\; \det A = 1\}
\]

元素可以写成：
\[
A = \begin{pmatrix} \alpha & -\bar{\beta} \\ \beta & \bar{\alpha} \end{pmatrix}, \quad |\alpha|^2 + |\beta|^2 = 1
\]

\begin{itemize}
  \item 维度：3
  \item 拓扑：$S^3$（三维球面）
  \item 单连通：$\pi_1(SU(2)) = 0$
  \item 同构于单位四元数群 $\text{Sp}(1)$
\end{itemize}

\textbf{关键关系：} $SU(2) \xrightarrow{2:1} SO(3)$（二重覆盖）

\textbf{物理意义：}
\begin{itemize}
  \item 电子自旋 $1/2$：旋转 $2\pi$ 后波函数变号，旋转 $4\pi$ 才回到自身
  \item 这就是为什么 $SU(2)$（而非 $SO(3)$）是自旋的正确对称群
\end{itemize}

\subsubsection{$SE(3)$：三维刚体运动}

\[
SE(3) = \left\{\begin{pmatrix} R & \mathbf{t} \\ \mathbf{0}^T & 1 \end{pmatrix} : R \in SO(3),\; \mathbf{t} \in \mathbb{R}^3\right\}
\]

\begin{itemize}
  \item 维度：6（3 旋转 + 3 平移）
  \item 拓扑：$SO(3) \times \mathbb{R}^3$（作为流形）
  \item 半直积结构：$SE(3) = SO(3) \ltimes \mathbb{R}^3$
\end{itemize}

\textbf{在 CG 中：} 刚体变换、骨骼动画、相机运动。

\subsection{李群的同态与覆盖}

\textbf{覆盖映射：} $\phi: \tilde{G} \to G$ 是局部同胚的满同态。

\textbf{万有覆盖群：} 单连通的覆盖群。

\begin{center}
\begin{tabular}{p{3cm}p{3cm}p{3cm}p{4cm}}
\toprule
\textbf{群} & \textbf{万有覆盖} & \textbf{覆盖映射} & \textbf{核} \\
\midrule
$SO(3)$ & $SU(2) \cong S^3$ & 2:1 & $\{\pm I\} \cong \mathbb{Z}_2$ \\
$SO(2) \cong S^1$ & $\mathbb{R}$ & $\theta \mapsto e^{i\theta}$ & $\mathbb{Z}$ \\
$SO(1,3)^+$（Lorentz） & $SL(2,\mathbb{C})$ & 2:1 & $\{\pm I\}$ \\
\bottomrule
\end{tabular}
\end{center}

%==============================================================
\section{李代数}
\label{sec:lie-algebra}
%==============================================================

\subsection{定义}

\textbf{李代数} $\mathfrak{g}$ 是向量空间，配备双线性运算 $[\cdot,\cdot]: \mathfrak{g} \times \mathfrak{g} \to \mathfrak{g}$（李括号），满足：
\begin{enumerate}
  \item 反对称：$[X, Y] = -[Y, X]$
  \item Jacobi 恒等式：$[X, [Y, Z]] + [Y, [Z, X]] + [Z, [X, Y]] = 0$
\end{enumerate}

\textbf{李群与李代数的关系：} 李代数 $\mathfrak{g} = T_e G$（李群在单位元处的切空间），李括号来自向量场的对易子。

\subsection{结构常数}

选基 $\{T_a\}$，$a = 1, \ldots, \dim\mathfrak{g}$：
\[
[T_a, T_b] = f_{ab}^{\;\;c}\, T_c
\]

$f_{ab}^{\;\;c}$ 是\textbf{结构常数}，完全决定了李代数的结构。

\textbf{Jacobi 恒等式的约束：}
\[
f_{ab}^{\;\;d}f_{dc}^{\;\;e} + f_{bc}^{\;\;d}f_{da}^{\;\;e} + f_{ca}^{\;\;d}f_{db}^{\;\;e} = 0
\]

\subsection{指数映射}

\[
\exp: \mathfrak{g} \to G, \quad X \mapsto e^X = \sum_{k=0}^{\infty} \frac{X^k}{k!}
\]

\textbf{性质：}
\begin{itemize}
  \item $\exp(0) = e$（单位元）
  \item $\exp(tX)$ 是过单位元的一参数子群
  \item 在单位元附近是局部微分同胚
  \item 一般不是满射（如 $SL(2,\mathbb{R})$）
\end{itemize}

\textbf{Baker--Campbell--Hausdorff 公式：}
\[
\log(e^X e^Y) = X + Y + \frac{1}{2}[X,Y] + \frac{1}{12}[X,[X,Y]] - \frac{1}{12}[Y,[X,Y]] + \cdots
\]

\textbf{意义：} 群乘法在李代数层面由李括号完全决定（局部）。

\subsection{经典李代数}

\subsubsection{$\mathfrak{so}(3)$}

\[
\mathfrak{so}(3) = \{A \in \mathbb{R}^{3\times 3} : A^T = -A\}
\]

基：
\[
J_1 = \begin{pmatrix} 0&0&0\\0&0&-1\\0&1&0 \end{pmatrix}, \quad
J_2 = \begin{pmatrix} 0&0&1\\0&0&0\\-1&0&0 \end{pmatrix}, \quad
J_3 = \begin{pmatrix} 0&-1&0\\1&0&0\\0&0&0 \end{pmatrix}
\]

对易关系：$[J_i, J_j] = \varepsilon_{ijk}J_k$

\textbf{与角动量的关系：} 量子力学中 $[\hat{L}_i, \hat{L}_j] = i\hbar\varepsilon_{ijk}\hat{L}_k$，正是 $\mathfrak{so}(3)$ 的结构（乘以 $i\hbar$）。

\subsubsection{$\mathfrak{su}(2)$}

\[
\mathfrak{su}(2) = \{A \in \mathbb{C}^{2\times 2} : A^\dagger = -A,\; \text{tr}\,A = 0\}
\]

基（Pauli 矩阵除以 $2i$）：
\[
\sigma_1 = \begin{pmatrix}0&1\\1&0\end{pmatrix}, \quad
\sigma_2 = \begin{pmatrix}0&-i\\i&0\end{pmatrix}, \quad
\sigma_3 = \begin{pmatrix}1&0\\0&-1\end{pmatrix}
\]

$T_a = \frac{1}{2}\sigma_a$，对易关系：$[T_a, T_b] = i\varepsilon_{abc}T_c$

\textbf{$\mathfrak{su}(2) \cong \mathfrak{so}(3)$}（作为李代数同构），但 $SU(2) \not\cong SO(3)$（作为群不同构，拓扑不同）。

\subsubsection{$\mathfrak{so}(1,3)$：Lorentz 代数}

\[
\mathfrak{so}(1,3) = \{A : A^T\eta + \eta A = 0\}, \quad \eta = \text{diag}(-1,1,1,1)
\]

6 个生成元：3 个旋转 $J_i$ + 3 个 boost $K_i$

对易关系：
\begin{align*}
[J_i, J_j] &= i\varepsilon_{ijk}J_k \\
[J_i, K_j] &= i\varepsilon_{ijk}K_k \\
[K_i, K_j] &= -i\varepsilon_{ijk}J_k
\end{align*}

\textbf{复化后分解：} $\mathfrak{so}(1,3)_\mathbb{C} \cong \mathfrak{su}(2)_\mathbb{C} \oplus \mathfrak{su}(2)_\mathbb{C}$

令 $A_i = \frac{1}{2}(J_i + iK_i)$，$B_i = \frac{1}{2}(J_i - iK_i)$，则 $[A_i, A_j] = i\varepsilon_{ijk}A_k$，$[B_i, B_j] = i\varepsilon_{ijk}B_k$，$[A_i, B_j] = 0$。

\textbf{物理意义：} 这解释了为什么粒子按 $(j_1, j_2)$ 分类（两个“自旋”量子数）：
\begin{itemize}
  \item $(0, 0)$：标量（Higgs）
  \item $(\frac{1}{2}, 0)$：左手 Weyl 旋量
  \item $(0, \frac{1}{2})$：右手 Weyl 旋量
  \item $(\frac{1}{2}, \frac{1}{2})$：向量（光子）
  \item $(1, 0)$ 或 $(0, 1)$：反对称张量
\end{itemize}

\subsubsection{$\mathfrak{su}(n)$}

\[
\mathfrak{su}(n) = \{A \in \mathbb{C}^{n\times n} : A^\dagger = -A,\; \text{tr}\,A = 0\}
\]

维度：$n^2 - 1$

\begin{itemize}
  \item $\mathfrak{su}(2)$：3 维，弱相互作用
  \item $\mathfrak{su}(3)$：8 维，强相互作用（量子色动力学）
\end{itemize}

\subsection{Killing 形式与半单性}

\textbf{Killing 形式：}
\[
B(X, Y) = \text{tr}(\text{ad}_X \circ \text{ad}_Y)
\]

其中 $(\text{ad}_X)(Y) = [X, Y]$。

\textbf{Cartan 判据：} $\mathfrak{g}$ 半单 $\iff$ $B$ 非退化。

\textbf{紧致性：} $\mathfrak{g}$ 是紧致李代数 $\iff$ $B$ 负定。

\subsection{伴随表示}

\textbf{伴随表示：} $\text{Ad}: G \to GL(\mathfrak{g})$，$\text{Ad}_g(X) = gXg^{-1}$

\textbf{无穷小版本：} $\text{ad}: \mathfrak{g} \to \mathfrak{gl}(\mathfrak{g})$，$\text{ad}_X(Y) = [X, Y]$

\textbf{在物理中：} 规范场（胶子、$W$ 玻色子）在伴随表示下变换。